\documentclass{handout}
\usepackage{handoutSetup}\usepackage{handoutShortcuts}

\provideboolean{CRRAUtil}\setboolean{CRRAUtil}{true}\newcommand{\ifCRRA}{\ifthenelse{\boolean{CRRAUtil}}}
\provideboolean{PermShks}\setboolean{PermShks}{true}\newcommand{\ifPerm}{\ifthenelse{\boolean{PermShks}}}

% If run from the command line script with option handout, set HandoutTF to true
\provideboolean{HandoutTF}\setboolean{HandoutTF}{false}\opt{handout}{\setboolean{HandoutTF}{true}}


\begin{document}
\handoutHeader

\begin{verbatimwrite}{\jobname.title}
Sticky Expectations and Consumption Dynamics
\end{verbatimwrite}

\handoutNameMake

Consider a consumer subject to the dynamic budget constraint
\begin{equation}\begin{gathered}\begin{aligned}
        \bLev_{t+1} & =  (\bLev_{t}+\yLev_{t}-\cLev_{t})\Rfree
\end{aligned}\end{gathered}\end{equation}
where $\bLev_{t}$ is beginning-of-period bank balances,
$\yLev_{t}$ is current labor income, and $\Rfree=(1+\rfree)$ is the constant
interest factor.  Actual labor income $\yLev$ is permanent labor income $\pLevBF$ modified by a 
transitory shock factor $\tShkEmp$:
\begin{equation}\begin{gathered}\begin{aligned}
  \yLev_{t+1} & =  \pLevBF_{t+1} \tShkEmp_{t+1}
\end{aligned}\end{gathered}\end{equation}
where $\Ex_{t}[\tShkEmp_{t+n}]=1~\forall~n>0$.  Permanent labor income grows by a 
predictable factor $\WGro$ from period to period:
\begin{equation}\begin{gathered}\begin{aligned}
  \pLevBF_{t+1} & =  \WGro \pLevBF_{t}\ifPerm{\psi_{t+1}}{},
\end{aligned}\end{gathered}\end{equation}
so that the expected present discounted value of permanent labor income (`human
wealth') for an infinite-horizon consumer is
\begin{equation}\begin{gathered}\begin{aligned}
  \hLev_{t} & =  \left(\frac{\pLevBF_{t}}{1-\WGro/\Rfree}\right).
\end{aligned}\end{gathered}\end{equation}

We will assume that the consumer behaves according to the consumption rule
\begin{equation}\begin{gathered}\begin{aligned}
  \cLev_{t} & =  \underbrace{(\bLev_{t}+(\tShkEmp_{t}-1)\pLevBF_{t}+\hLev_{t})}_{\equiv \wAllLev_{t}}\underbrace{(\rfree/\Rfree)}_{\equiv \MPC}, \label{eq:DeatonPIH}
\end{aligned}\end{gathered}\end{equation}
where $\MPC$ is the `marginal propensity to consume' out of total wealth $\wAllLev$.\footnote{This is the optimal consumption function for a utility-maximizing consumer with $\Rfree \Discount = 1$ if that consumer has quadratic utility (\cite{hallRandomWalk}) or if the consumer has CRRA utility and perfect foresight and anticipates $\tShkEmp_{t+n}=\pShk_{t+n}=1~\forall~n>0$.  See \handoutC{ConsumptionFunction} for a derivation of this consumption function under quadratic utility, and \handoutC{PerfForesightCRRA} for the derivation in the perfect foresight CRRA case.  \cite{deatonUnderstandingC} argues that the `Permanent Income Hypothesis' should be {\bf defined} as the hypothesis that consumption is determined according to \eqref{eq:DeatonPIH}; but this differs sharply from \cite{friedmanATheory}'s definition, and has not become universally accepted.}

Under these circumstances, \handoutC{RandomWalk} shows that 
consumption will follow a random walk,
\begin{equation}\begin{gathered}\begin{aligned}
\Delta \cLev_{t+1} & =  \error_{t+1}, \label{eq:rwc}
\\ \Ex_{t}[\error_{t+n}] & =  0 ~\forall~n>0.
\end{aligned}\end{gathered}\end{equation}

Now assume that the economy is populated by a set of measure one of
consumers indexed by a superscript $i$ distributed uniformly along the unit
interval.  Per capita values of all variables, designated by the upper
case, are the integral over all individuals in the economy, 
as in \Aggregation, so that 
\begin{equation}\begin{gathered}\begin{aligned}
        \CLev_{t}  & =  \int_{0}^{1} \cLev_{t}^{i} di
\\       & =  \int_{0}^{1} (\rfree/\Rfree)  \wAllLev_{t}^{i} di
\\   & =  \WAllLev_{t} \MPC
.
\end{aligned}\end{gathered}\end{equation}
This equation implies that an aggregate version of equation~\eqref{eq:rwc} holds,\footnote{The crucial feature of the model that allows us to aggregate 
analytically is the linearity of the consumption rule in $p, b, $ and 
$\TShkEmp$.}
  \begin{equation}\begin{gathered}\begin{aligned}
    \Delta \CLev_{t+1} & =  \error_{t+1}
.
  \end{aligned}\end{gathered}\end{equation}

In principle, we could allow each individual in this economy to experience a different transitory \ifPerm{and permanent}{} shock from every other individual in each period.  However, for our purposes
it is useful to assume that everyone experiences the same shocks in a given period; that is $\tShkEmp_{t}^{i}=\TShkEmp_{t}~\forall~t$\ifPerm{ and $\psi_{t}^{i} = \Psi_{t}~\forall~t$}{}.

Assuming (here and henceforth) that the growth factor for permanent income is $\WGro=1$, figure~\ref{fig:aftershock} shows the path of consumption (the solid
dots) for an economy populated by omniscient consumers who in periods
$t-n$ for $n>0$ had experienced
$\TShkEmp_{t-n}\ifPerm{=\Psi_{t-n}}{}=1$; that is, this economy has had
no shocks to income in the past.  (For convenience, the consumer is
assumed to have arrived in period $t$ with $\BLev_{t}=0$).  In period
$t$ the consumer draws $\TShkEmp_{t}=2$ \ifPerm{and $\Psi_{t}=1$}{};
thereafter $\ifPerm{\Psi_{t+n}=}{}\TShkEmp_{t+n}=1$.  The figure shows
$\CLev_{t-2},\CLev_{t-1},\CLev_{t},\Ex_{t}[\CLev_{t+1}],\Ex_{t}[\CLev_{t+2}],\ldots$.
Figure~\ref{fig:epidshock} similarly shows the path of $\BLev$, again
as black dots.

\begin{figure}[h]
\begin{center}
 \caption{Path of $\CLev$ after a shock $\TShkEmp_{t}=2$, Omniscient Consumers}
 \label{fig:aftershock}
\includegraphics[width=6in]{../Figures/cPlot}
 \caption{Path of $\BLev$ after a shock $\TShkEmp_{t}=2$, Omniscient Consumers}
 \label{fig:epidshock}
\includegraphics[width=6in]{../Figures/bPlot}
\end{center}
\end{figure}

\begin{comment}
\renewcommand{\bLev}{\ensuremath{b}}
\renewcommand{\cLev}{\ensuremath{c}}
\renewcommand{\yLev}{\ensuremath{y}}
\renewcommand{\hLev}{\ensuremath{\mathsf{h}}}
\renewcommand{\zLev}{\ensuremath{z}}
\renewcommand{\error}{\epsilon}
\renewcommand{\BLev}{\ensuremath{B}}
\renewcommand{\CLev}{\ensuremath{C}}
\renewcommand{\YLev}{\ensuremath{Y}}
\renewcommand{\ZLev}{\ensuremath{Z}}
\renewcommand{\Err}{\epsilon}

A convenient feature of this model is that the problem can be normalized by the level of permanent income.
Using nonbold variables to represent the normalized version of bold variables, and for convenience defining $\hat{\TShkEmp}_{t}=\TShkEmp_{t}-1$ the consumption rule 
becomes
\begin{equation}\begin{gathered}\begin{aligned}
  \cLev_{t} & =    (\bLev_{t}+\hat{\TShkEmp}_{t} + \hLev_{t}) \MPC \label{eq:cratio}
.
\end{aligned}\end{gathered}\end{equation}
\end{comment}

Now suppose that not every consumer updates expectations in every 
period.  Instead, expectations are `sticky': each consumer 
updates with probability $\Pi$ in each period.  
Whether the consumer at location $i$ updates in period $t$ is determined
by the realization of the dichotomous random variable
\begin{equation*}
\pi^{i}_{t} = % PiInd in the software
\begin{cases}
   1 & \text{if consumer $i$ updates in period $t$}
\\ 0 & \text{if consumer $i$ does not update in period $t$}  ,
\end{cases}
\end{equation*}
and each period's updaters are chosen randomly
such that a constant proportion $\Pi$ update in each period:
\begin{equation*}\begin{gathered}\begin{aligned}
   \Ex_{t}^{i}[\pi^{i}_{t+1}] &=  \Pi~\forall~t
\\ \int_{0}^{1} \pi^{i}_{\tau} di & =  \Pi ~\forall~\tau
.
\end{aligned}\end{gathered}\end{equation*}

It will also be convenient to define the date of consumer
$i$'s most recent update; we call this object $\tau^{i}_{t}$.  
\begin{comment}
We designate the value of a variable $\bullet$ at date $t$ for a
consumer who last updated his expectations at date $t-n$ by
$\bullet_{t|=t-n}$ where the part of the subscript following $|$ indicates criteria
that must be matched by the consumer's (or group's) $\tau$.  
\end{comment}

\begin{comment}
For
example, consider a consumer who updated his expectations most
recently in period $t-3$ so that
$n(i)=3$.  % Designating $\{\cLev^{\tau}_{\tau},\cLev_{\tau|\tau}\}$ as the level of consumption a fully-informed (omniscient) consumer with completely up-to-date expectations would choose, t
This consumer's actual consumption in period $t$ will be
\begin{equation}\begin{gathered}\begin{aligned}
%        \cLev^{t-n(i)}_{t} & =  \cLev^{t-3}_{t-3}. \\
      \cLev^{i}_{t} & =  \cLev_{t-3|=t-3}.
\end{aligned}\end{gathered}\end{equation}
\end{comment}

\begin{comment}
We are assuming that the probability of adjusting one's expectations
is independent of the level of income or wealth; therefore, the
average levels of wealth and income among those consumers who adjust
their expectations this period will be the same as the average levels
of wealth and income in the economy as a whole.
\end{comment}

\begin{comment}
Dropping the $i$ superscript for notational simplicity, the level of consumption
per capita that would prevail if all consumers were to update in
period $t$ is
\begin{equation}\begin{gathered}\begin{aligned}
        \CLev_{t|=t} & =  \int_{0}^{1} \cLev_{t|=t} di.
\end{aligned}\end{gathered}\end{equation}
But since the set of consumers who updated was randomly selected from
the population, the average level of consumption-per-capita for the
updaters must equal the average level of consumption-per-capita that
would characterize the population as a whole if everyone in the
economy were to update.
Thus, an alternative measure of the same thing is given by
\end{comment}

We need a notation to represent sets of consumers defined by the period of their most recent
update.  We denote such a set by the condition
on $\tau^{i}_{t}$; for example, the set of consumers whose most recent update, as of
date $t$, was prior to period $t-1$ would be $\mathcal{T} = \{\tau^{i}_{t} < t-1\}$.  
We denote the per-capita value of a variable $\bullet$, among consumers in a set $\mathcal{T}$
as of date $t$, by $\bullet_{t|\mathcal{T}}$.  Dropping the $i$ superscripts to reduce clutter, per-capita consumption among households who have updated in period $t$ is therefore
\begin{equation}\begin{gathered}\begin{aligned}
        \CLev_{t|\tau_{t}=t} & =  \Pi^{-1} \int_{0}^{1} \pi_{t} \cLev_{t} di
.
\end{aligned}\end{gathered}\end{equation}

In periods when expectations are not updated, the
consumer continues to spend the same amount as in the most
recent period when his expectations {\bf were} updated.\footnote{This makes sense because under the \cite{hallRandomWalk} assumptions the expected change in consumption is zero under our assumption that $\Rfree \Discount=1$.}  If the economy is large the proportion of consumers who update their 
expectations every period will be $\Pi$.\footnote{For consumers who are not updating at date $t$, $\cLev_{t} > 0$ but $\pi_{t} = 0$, so the integral produces the sum of consumption among only those $i$ who are updating in period $t$.}  Average consumption among those who are not updating in the current 
period (for whom $1-\pi_{t}=1$) is then 
\begin{equation}\begin{gathered}\begin{aligned}
  \CLev_{t|\tau_{t} < t} & =  (1-\Pi)^{-1} \int_{0}^{1} (1-\pi_{t}) \cLev_{t} di
\\ & =  \CLev_{t-1}
\end{aligned}\end{gathered}\end{equation}
because consumption per capita among those who are not updating in the
current period is (by assumption) identical to their consumption per
capita in the prior period, which must match aggregate consumption per
capita in the prior period because the set who do {\bf not} update
today is randomly selected from the $t-1$ population.

Now note that 
\begin{equation}\begin{gathered}\begin{aligned}
          \CLev_{t+1} & =  \Pi \CLev_{t+1|\tau_{t+1}=t+1} + (1-\Pi) \overbrace{\CLev_{t+1|\tau_{t+1}<t+1}}^{=\CLev_{t}}
\\ \Delta \CLev_{t+1} & =  \Pi (\CLev_{t+1|\tau_{t+1}=t+1}-\CLev_{t|\tau_{t}=t}) + (1-\Pi) \Delta \CLev_{t} \label{eq:DCtp1}
\end{aligned}\end{gathered}\end{equation}
and 
\begin{equation}\begin{gathered}\begin{aligned}
          \CLev_{t} & =  \Pi \CLev_{t|\tau_{t}=t} + (1-\Pi) \CLev_{t-1}
\\ \Delta \CLev_{t} & =  \Pi (\CLev_{t|\tau_{t}=t}-\CLev_{t-1}) 
\\  & =  \Pi (\CLev_{t|\tau_{t}=t}-\CLev_{t}+\CLev_{t}-\CLev_{t-1})
\\ (1-\Pi) \Delta \CLev_{t} & =  \Pi (\CLev_{t|\tau_{t}=t}-\CLev_{t}) \label{eq:DCvsCmC}
\end{aligned}\end{gathered}\end{equation}
while, defining $\hat{\TShkEmp}=\TShkEmp-1$ and $\hat{\Psi}=\Psi-1$,
\begin{equation}\begin{gathered}\begin{aligned}
  \CLev_{t+1|\tau_{t+1}=t+1} & =  \left(\overbrace{(\BLev_{t}+\hat{\TShkEmp}_{t}-\CLev_{t})\Rfree}^{=\BLev_{t+1}}+\hat{\TShkEmp}_{t+1}\right)\MPC+1\ifPerm{+\hat{\Psi}_{t+1}}{}
\\  \CLev_{t+1|\tau_{t+1}=t+1} & =  \left((\BLev_{t}+\hat{\TShkEmp}_{t}-\CLev_{t}+\CLev_{t|\tau_{t}=t}-\CLev_{t|\tau_{t}=t})\Rfree+\hat{\TShkEmp}_{t+1}\right)\MPC+1\ifPerm{+\hat{\Psi}_{t+1}}{}
\\ & =  \underbrace{\left((\BLev_{t}+\hat{\TShkEmp}_{t}-\CLev_{t|\tau_{t}=t})\Rfree+\hat{\TShkEmp}_{t+1}\right)\MPC+1\ifPerm{+\hat{\Psi}_{t+1}}{}}_{=\Ex_{t}[\CLev_{t+1|\tau_{t}=t}]+\hat{\TShkEmp}_{t+1}\MPC\ifPerm{+\hat{\Psi}_{t+1}}{}}+(\CLev_{t|\tau_{t}=t}-\CLev_{t})\Rfree\MPC \nonumber
\\ & =  \CLev_{t|\tau_{t}=t}+\hat{\TShkEmp}_{t+1}\MPC\ifPerm{+\hat{\Psi}_{t+1}}{}+(\CLev_{t|\tau_{t}=t}-\CLev_{t})\MPC \Rfree \label{eq:DCtp1Gtp1}
%\\ \CLev_{t+1|\tau_{t+1}=t+1}-\CLev_{t|\tau_{t}=t} & =  \hat{\TShkEmp}_{t+1}\MPC\ifPerm{+\hat{\Psi}_{t+1}}{}+(\CLev_{t|\tau_{t}=t}-\CLev_{t})\MPC \Rfree 
\end{aligned}\end{gathered}\end{equation}
where \eqref{eq:DCtp1Gtp1} follows its predecessor since, among consumers who have updated in period $t$, the random walk proposition 
says that $\Ex_{t}[\CLev_{t+1|\tau_{t}=t}]=\CLev_{t,\tau_{t}=t}$.  Subtracting $\CLev_{t|\tau_{t}=t}$ from both sides of \eqref{eq:DCtp1Gtp1} and substituting the result into \eqref{eq:DCtp1} yields
\begin{equation}\begin{gathered}\begin{aligned}
 \Delta \CLev_{t+1} & =  (1-\Pi) \Delta \CLev_{t}+
\underbrace{\Pi (\CLev_{t|\tau_{t}=t}-\CLev_{t})}_{=(1-\Pi)\Delta \CLev_{t}~\text{from \eqref{eq:DCvsCmC}}}\underbrace{\Rfree\MPC}_{=\rfree} + \underbrace{\Pi\left(\hat{\TShkEmp}_{t+1}\MPC\ifPerm{+\hat{\Psi}_{t+1}}{}\right)}_{\equiv \xi_{t+1}} 
\\ & =  (1-\Pi) \Delta \CLev_{t}+ (1-\Pi) \rfree \Delta \CLev_{t} + \xi_{t+1}
\\ & =  (1-\Pi) \Rfree \Delta \CLev_{t} + \xi_{t+1}
 \end{aligned}\end{gathered}\end{equation}
where $\xi$ is a white noise variable ($\Ex_{t}[\xi_{t+n}]=0~\forall~n>0$).

We are finally in position to show how aggregate consumption and wealth
would respond in this economy to a transitory positive shock to
aggregate labor income like the one considered above for the
omniscient model.

Consider the case of a positive shock of size $\hat{\TShkEmp}_{t}=1$, as before.  In the first period 
consumption rises only by $\Pi \MPC$, rather than the full amount 
corresponding to the permanent income associated with the new level of 
wealth.  Therefore aggregate wealth in period $t+1$ will be greater 
than it would have been in the omniscient model.  Similarly for 
all subsequent periods.  Thus, in contrast with the omniscient 
model, the sluggish adjustment of consumption to the shock means that 
the shock has a permanent effect on the level of aggregate wealth, and 
therefore on the level of aggregate consumption.  (Figures~\ref{fig:StickyExpectationsCc} and \ref{fig:StickyExpectationsCb} depict the results).

\begin{figure}
\begin{center}
 \caption{Path of $\CLev$ after a shock $\TShkEmp_{t}=2$; Sticky Expectations in Red/Gray}
<\includegraphics[width=6in]{../Figures/StickyExpectationsCc}\label{fig:StickyExpectationsCc}
 \caption{Path of $\BLev$ after a shock $\TShkEmp_{t}=2$; Sticky Expectations in Red/Gray}
\includegraphics[width=6in]{../Figures/StickyExpectationsCb}\label{fig:StickyExpectationsCb}
\end{center}
\end{figure}

The sticky expectations model says that consumption growth today 
can be statistically related to any variable that is related to lagged 
consumption growth.  In particular, if lagged consumption growth is 
related to lagged income growth (as it certainly will be), then there 
{\bf should} be a statistically significant effect of lagged income 
growth on current consumption growth if expectations are sticky.

If the model derived here could be taken literally, it would suggest estimating
an equation of the form
\begin{equation}\begin{gathered}\begin{aligned}
        \Delta \CLev_{t+1} & =  \alpha_{0}+\alpha_{1} \Delta \CLev_{t} + \varepsilon_{t+1} \label{eq:eststick}
\end{aligned}\end{gathered}\end{equation}
and interpreting the coefficient $\alpha_{1}$ as a measure of 
$\Rfree (1-\Pi)$.

 


However, if there is potential measurement error in $\CLev_{t}$ the
coefficient obtained from estimating \eqref{eq:eststick} would be
biased toward zero for standard errors-in-variables reasons (just as
regressing consumption on actual income yields a downward-biased
estimate of the response of consumption to permanent income), which
means that the estimate of $\Pi$ would be biased toward 1 (i.e. the
omniscient model in which everyone adjusts all the time).  Under these
circumstances, direct estimation of \eqref{eq:eststick} would not be a
reliable way to estimate $\Pi$.

For estimation methods that get around this problem see \cite{som07}, \cite{cssIntlStickyC}, \cite{cosHousingWealth}.  Those papers consistently find that the proportion of updaters is about $\Pi=0.25$ per quarter, so that the serial correlation of `true' consumption growth is about $(1-\Pi)=0.75$ per quarter.

\input handoutBibMake.tex

%\ifthenelse{\boolean{HandoutTF}}{}{
\pagebreak\vfill\eject
\appendix
%}{} % End HandoutTF


\centerline{\Large Appendix}\medskip\medskip

This appendix provides additional derivations and notation useful for simulating the model.  
One way of interpreting consumers' behavior in this model is to attribute to them the beliefs
that would rationalize their actions.  Define $\WAllLev$ as the level of wealth (human and nonhuman)
that the consumer perceives.  Then the Deaton definition of the permanent income hypothesis is that 
\begin{equation}\begin{gathered}\begin{aligned}
  \CLevBF_{t} & =  \WAllLev_{t} \MPC
\end{aligned}\end{gathered}\end{equation}
and the reason consumption follows a random walk is that $\MPC=(\rfree/\Rfree)$ is precisely the amount
that ensures that $\Ex_{t}[\WAllLev_{t+1}] = \WAllLev_{t}$.

Writing the ``believed'' level of wealth as $\bar{\WAllLev}$, we could then interpret the
failure of the sticky expectations consumer to change his consumption during the period 
of nonupdating as reflecting his optimal forecast that, in the absence of further information,
$\bar{\WAllLev}_{t+1} = \bar{\WAllLev}_{t+2} + ... = \WAllLev_{t}$.  

To pursue this interpretation, it is useful to write the budget constraint more explicitly, 
as before; start with the constraint in levels, then decompose variables into ratios to permanent income (nonbold variables) and the level of permanent income:
\begin{equation}\begin{gathered}\begin{aligned}
    \BLevBF_{t+1} & =  \ALevBF_{t} \Rfree_{t+1}
\\  \BRat_{t+1}P_{t+1} & =  \ARat_{t}P_{t} \Rfree_{t+1}
\\  \BRat_{t+1} & =  \ARat_{t} \underbrace{(\Rfree_{t+1}/\WGro_{t+1})}_{\equiv \Rnorm_{t+1}}
\end{aligned}\end{gathered}\end{equation}
where we permit a time subscript on $\Rfree$ and $\WGro$ because we want to allow for the 
possibility that beliefs about the interest rate or growth rate might change over time.
\ifPerm{}{We also allow henceforth for the existence of permanent shocks to income, $\Psi_{t}$.}

Consider an economy that comes into existence in period $0$ with a population of 
consumers who are identical in every respect, including their beliefs about current and future
values of the economy's variables.  

First we examine the case where neither $\Rfree$ nor $\WGro$ can change after date $0$.
In that case, we can track the dynamics of believed and actual variables as follows.
\begin{equation}\begin{gathered}\begin{aligned}
  \bar{\BLevBF}_{t+1} & =  \Pi \BLevBF_{t+1} + (1-\Pi) \bar{\ALevBF}_{t} \Rfree
\\ \bar{\BRat}_{t+1}\bar{P}_{t+1} & =  \Pi \BRat_{t+1}P_{t+1} + (1-\Pi) \bar{\ARat}_{t}\bar{P}_{t} \Rfree
\\ \bar{\BRat}_{t+1} & =  \Pi \BRat_{t+1}\underbrace{(P_{t+1}/\bar{P}_{t+1})}_{\equiv Q_{t+1}} + (1-\Pi)\bar{\ARat}_{t} \underbrace{(\Rfree/\WGro \bar{\hat{\Psi}}_{t+1})}_{\equiv \bar{\Rnorm}_{t+1}} 
\end{aligned}\end{gathered}\end{equation}
where capturing the dynamics of the ratio of true permanent income to
believed permanent income $Q_{t+1}$ requires us to compute
\begin{equation}\begin{gathered}\begin{aligned}
  \bar{P}_{t+1} & =  \left(\Pi \Psi_{t+1} P_{t} + (1-\Pi) \bar{P}_{t}\right)\WGro
\\ \bar{\Psi}_{t+1} & \equiv  \bar{P}_{t+1}/(\bar{P}_{t} \WGro)
\\ Q_{t+1} & =  (\Psi_{t+1}/\bar{\Psi}_{t+1}) Q_{t}
\end{aligned}\end{gathered}\end{equation}
and so
\begin{equation}\begin{gathered}\begin{aligned}
  (\bar{\TShkEmp}_{t+1}-1)\bar{P}_{t+1} & =  \Pi (\TShkEmp_{t+1}-1) P_{t+1} + (1-\Pi) (1-1) \bar{P}_{t} \WGro
\\  \bar{\TShkEmp}_{t+1}-1 & =  \Pi (\TShkEmp_{t+1}-1) Q_{t+1}
%\begin{equation}\begin{gathered}\begin{aligned}
%  \bar{\YRat}_{t+1}\bar{P}_{t+1} & =  \Pi \TShkEmp_{t+1}P_{t+1} + (1-\Pi)\bar{P}_{t}\WGro
%\\ \bar{\YRat}_{t+1} & =  \Pi \TShkEmp_{t+1}Q_{t+1} + (1-\Pi)/\bar{\Psi}_{t+1}
\end{aligned}\end{gathered}\end{equation}
with the crucially useful fact that since by assumption neither $\Rfree$ nor $\WGro$ is 
changing, normalized human wealth does not change from
\begin{equation}\begin{gathered}\begin{aligned}
  \hRat & =  \left(\frac{1}{1-\WGro/\Rfree}\right)
\end{aligned}\end{gathered}\end{equation}
so that 
\begin{equation}\begin{gathered}\begin{aligned}
  \bar{\HLev}_{t} & =  \hRat \bar{P}_{t}
\\ \HLev_{t} & =  \hRat P_{t}
\\ \bar{\HLev}_{t+1} & =  \left(\Pi Q_{t+1} + (1-\Pi)/\bar{\Psi}_{t+1} \right) \hRat \bar{P}_{t+1}
\end{aligned}\end{gathered}\end{equation}
so that perceived wealth and consumption will be
\begin{equation}\begin{gathered}\begin{aligned}
  \bar{\WAllRat}_{t+1} & =  \bar{\BRat}_{t+1}+(\bar{\TShkEmp}_{t+1}-1)+\bar{\HRat}_{t+1}
\\ \bar{\CRat}_{t+1} & =  \bar{\WAllRat}_{t+1} \MPC
\\ \bar{\ARat}_{t+1} & =  \bar{\BRat}_{t+1}+\bar{\TShkEmp}_{t+1}-\bar{\CRat}_{t+1}
.
\end{aligned}\end{gathered}\end{equation}

Matters are more complex if expectations about $\Rfree$ and $\WGro$ are allowed to change
over time.

Suppose again that we begin our economy in period $0$ with population with homogeneous views: Everyone believes $\Rfree=\Rfree^{0}$ and $\WGro=\WGro^{0}$; so long as these views are universally held in the population, aggregate dynamics are captured by the foregoing analysis.  

Suppose, however, that in some period $n>0$ the economy's `true' values of $\Rfree$ or $\WGro$ change.  
Updating consumers see this change immediately.  But nonupdaters will not discover the changed nature
of the economy's dynamics until they update again.  

We capture this modification to the model by keeping track of the
aggregate values of the variables for the set of consumers who adhere
to each differing opinion, along with the population mass associated
with the different opinions.  Specifically, suppose there are $J$
different opinions in the population, each of whom constitutes
population mass $\mathcal{L}^{j}_{t}$ such that $\sum_{j}
\mathcal{L}^{j}_{t} = 1~\forall~t$.  Then for each such population, it will be
necessary to keep track of their average beliefs about macroeconomic
variables.  

Suppose, for example, that through period $x$ there have been only $j-1$ different 
opinion groups in the population.  In period $x+1$ either $\WGro$ or $\Rfree$ changes.  
We need then to define group $j$ by $\Rfree^{j}=\Rfree_{x}$ and $\WGro^{j}=\WGro_{x}$ and
to define $\bar{\ARat}^{j}_{x} \equiv \ARat_{x}$, $\bar{P}^{j}_{x}=P_{x}$, and so on.
We will henceforth need to keep track of dynamics of the consumers who remain in belief group
$j$ by, e.g.,
\begin{equation}\begin{gathered}\begin{aligned}
  \bar{\BLevBF}^{j}_{t+1} & =  \ALevBF^{j}_{t} \Rfree^{j}
\\ \bar{\BRat}^{j}_{t+1}\bar{P}^{j}_{t+1} & =  \bar{\ARat}^{j}_{t}\bar{P}^{j}_{t} \Rfree^{j}
\\ \bar{\BRat}^{j}_{t+1} & =  \bar{\ARat}^{j}_{t} \underbrace{(\Rfree^{j}/\WGro^{j} \bar{\Psi}^{j}_{t+1})}^{j}_{\equiv \bar{\Rnorm}^{j}} 
\\ \bar{P}^{j}_{t+1} & =  \bar{P}^{j} \WGro^{j}
\end{aligned}\end{gathered}\end{equation}
while we need to keep track of the populations of the differing groups by, e.g., 
\begin{equation}\begin{gathered}\begin{aligned}
    \mathcal{L}^{j}_{t} & =  \Pi
\\  \mathcal{L}^{j}_{t+1} & =  \Pi(1-\Pi)
\\  \mathcal{L}^{j}_{t+2} & =  \Pi(1-\Pi)^{2}
\end{aligned}\end{gathered}\end{equation}
and so on.  These population dynamics continue forever, but the population of households continuing
to hold any specific belief configuration dwindles toward zero as time progresses.  

Aggregate variables for the population as a whole can be constructed as the population-weighted
sums across all the differing belief groups, weighted by their masses:
\begin{equation}\begin{gathered}\begin{aligned}
  \bar{P}_{t} & =  \Pi P_{t} + (1-\Pi) \sum_{j} \mathcal{L}^{j}_{t} \bar{P}^{j}_{t}  
\end{aligned}\end{gathered}\end{equation}
and note that if beliefs change back to a configuration that has been seen before it
is possible to add the population mass and aggregate values of the variables associated
with the new population with that belief configuration to the corresponding figures for 
the old population that holds the same beliefs.  This reduces the number of groups that
the simulations must track in the case where beliefs switch between a limited number of 
distinct values.  



\end{document}


